\section{恒瑞医药：收入、利润与SOTP闭环}

恒瑞2025年收入316.29亿元，其中创新药销售163.42亿元、对外许可33.92亿元，成熟与其他业务代理约118.95亿元。2026Q1创新药销售45.26亿元、许可收入7.87亿元、归母净利22.82亿元。模型不再把全公司统一套创新药倍数，而是分别估算创新药、成熟业务与许可价值。

\begin{exhibitbox}[Exhibit 12：恒瑞收入—利润—SOTP桥]
\begin{tabularx}{\textwidth}{L{1.4cm}R{1.35cm}R{1.35cm}R{1.35cm}R{1.25cm}R{1.25cm}R{1.25cm}X}
\toprule
\textbf{情景} & \textbf{创新销售} & \textbf{许可收入} & \textbf{成熟其他} & \textbf{总收入} & \textbf{净利润} & \textbf{EPS} & \textbf{SOTP/股} \\
\midrule
熊 & 196.1 & 20.0 & 109.4 & 325.5 & 76.5 & 1.15 & 48.97 \\
基准 & 212.4 & 31.7 & 113.0 & 357.1 & 95.4 & 1.44 & 68.54 \\
牛 & 220.6 & 40.0 & 119.0 & 379.6 & 110.1 & 1.66 & 79.19 \\
\bottomrule
\end{tabularx}
\sourcenote{data/growth\_driver\_model.json与data/hengrui\_sotp\_model\_20260712.json。}
\end{exhibitbox}

华泰2026-03-26原始PDF的未折价SOTP为5918亿元：创新药DCF 5408亿元、仿制药及其他192亿元、对外授权317亿元，WACC 7.4\%、永续增长2.5\%，对应89.16元。AStock牛市SOTP仍折价至5256亿元／79.19元，华泰89.16元只作为独立10\%Street锚，避免在基本面情景与外部锚中重复计价。

现价55.75元若按40倍PE，隐含EPS约1.39元、净利润约92.5亿元，已计入大部分基准盈利但未计入完整SOTP。已核验公司Q1证据、华源2026-05-07、西南2026-07-01与华泰2026-03-25原始PDF；组合层面已有创新药收入、BD确认、NDA／III期进展和分部SOTP闭环。产品级峰值销售、分阶段PoS与未来BD确认日属于正式披露边界，因此AStock对创新药DCF折价并对许可收入概率加权，不向任何单品补入未经披露的峰值销售。

\section{工业富联：H1预告到H2利润门槛}

工业富联H1归母净利中值239亿元、扣非232亿元；Q1归母105.95亿元，隐含Q2约133.05亿元。公司披露AI服务器收入同比增长超过230\%、800G以上交换机出货增长140\%，构成比泛AI叙事更直接的经营证据。

\begin{exhibitbox}[Exhibit 13：工业富联盈利敏感度]
\begin{tabularx}{\textwidth}{L{1.4cm}R{1.55cm}R{1.55cm}R{1.25cm}R{1.15cm}R{1.35cm}X}
\toprule
\textbf{情景} & \textbf{收入} & \textbf{净利润} & \textbf{EPS} & \textbf{PE} & \textbf{价值} & \textbf{验证} \\
\midrule
熊 & 13000.0 & 500.6 & 2.52 & 20x & 50.40 & 平台放缓/毛利承压 \\
基准 & 14803.7 & 560.1 & 2.82 & 27x & 76.14 & H2净利321.1亿元 \\
牛 & 16000.0 & 646.0 & 3.25 & 30x & 97.50 & 平台与利润超预期 \\
\bottomrule
\end{tabularx}
\end{exhibitbox}

现价66.27元对应基准EPS约23.5倍PE；若按27倍基准倍数，现价隐含净利润约488亿元，低于基准560.1亿元，但安全垫不厚。正式目标74.06元要求H2净利润至少321.1亿元、AI服务器增速保持100\%以上、毛利率不低于7\%。已核验官方H1预告、金元原始PDF、华泰平台报告和预告后完整研报页：除AI服务器收入+230\%、800G出货+140\%外，公司还披露主要客户份额提升、联合研发的下一代产品H2量产；CPO样机生产、上万台指引与Rubin爬坡提供代理。客户名称分配、单品ASP、精确利用率与AI分部毛利属于正式边界，因此只用合并PE，不给AI分部单独高倍数。

\section{四只高增长核心的证据闭环}

高增长估值必须同时连接收入代理、毛利、费用、净利润、EPS和当前价格隐含门槛。以下四只均已完成来源检查与代理证据闭环；差别不是“有没有证据”，而是证据能支持合并估值、分部SOTP，还是只能支持降权后的条件情景。

\fontsize{7.5pt}{10pt}\selectfont
\noindent
\begin{tabularx}{\textwidth}{L{1.35cm}L{5.10cm}L{4.00cm}X}
\toprule
\textbf{公司} & \textbf{直接／代理证据} & \textbf{正式披露边界} & \textbf{估值后果} \\
\midrule
恒瑞医药 & 创新药收入、BD确认、NDA／III期进展、12笔BD与组合SOTP & 单品峰值、阶段PoS、未来BD确认日 & 创新药DCF折价，许可收入概率加权 \\
工业富联 & AI服务器+230\%、800G+140\%、大客户联合研发、CPO／Rubin代理 & 客户分配、单品ASP、利用率、AI分部毛利 & 只用合并PE，不给AI分部额外倍数 \\
中兴通讯 & 算力收入+150\%、占比24.6\%，服务器存储+200\%、Q1占比27\%；头部互联网、运营商、政务金融客户类别；国金2026E收入1728.08亿元、净利73.73亿元、EPS 1.541元 & 具名客户、设备量、合同ASP、算力分部毛利 & 仅合并PE；概率值低于现价，条件观察 \\
江波龙 & LTA／MOU、自研SPU／HLC、自有封测、AMD联合调优；NAND／DRAM价格、平台认证、企业级存储及客户行业代理 & 合同价、库存成本层、H2具名客户订单额 & 11x／14x／17x周期PE与35\%熊市概率 \\
\bottomrule
\end{tabularx}
\normalsize
\sourcenote{data/growth\_driver\_model.json；每行均保存检查来源、正式边界和估值后果。}
